\subsection{Algebraic extensions of ordered fields}

Let K be an ordered field, and f a polynomial in K{[}X{]}. We will say that f changes sign in K if there exist two elements a and b in K such that \(f(a)\) \(f(b) < 0\) ; then we say that f changes sign between a and b.

PROPOSITION 3. --- Let K be an ordered field and f an irreducible polynomial over K which changes sign between a and b in K. Then the extension \(E = K[X]/(f)\) of K admits the structure of an ordered extension.

We will argue by induction on the degree n of f.~For n = 1 the proof is trivial. Suppose the result holds for degrees \(\leqslant n - 1\) , and let us prove it for degree n by contradiction; by Th. 1 we are thus assuming a relation of the form

\(1 + \sum_{i}p_{i}f_{i}^{2}(\mathbf{X})\equiv 0 (\mathrm{mod} f(\mathbf{X}))\) , where \(f_{i}\in \mathbf{K}[\mathbf{X}], p_{i}\in \mathbf{K}\) and \(p_i\geqslant 0\)

Without loss of generality we may suppose that the \(f_{i}\) have degrees \(\leqslant n - 1\) (IV, p.~11, Cor). Then

\[
1 + \sum_ {i} p _ {i} f _ {i} ^ {2} (\mathbf {X}) = h (\mathbf {X}) f (\mathbf {X})
\]

where \(h \neq 0\) has degree at most n - 2. Replacing X by a and b in the above inequality, we see that \(h(a) f(a) > 0\) and \(h(b) f(b) > 0\) . Since f changes sign between a and b by hypothesis, we conclude that \(h(a) h(b) < 0\) . Then we have a similar inequality for one of the irreducible factors \(g(\mathbf{X})\) of \(h(\mathbf{X})\) : that is \(g(a)g(b) < 0\) . But \(1 + \sum_{i}p_{i}f_{i}^{2}(\mathbf{X})\equiv 0 (\mathrm{mod} g(\mathbf{X}))\) , which shows that the field \(\mathrm{K}[\mathrm{X}] / (g)\) cannot be an ordered extension of \(\mathbf{K}\) (Th. 1), contrary to the induction hypothesis.

Remark. --- There exist irreducible polynomials f over an ordered field K which do not change sign in K, but such that \(\mathrm{K}[\mathrm{X}]/(f)\) admits the structure of an ordered extension of K (cf.~VI, p.~43, Ex. 26, c)).

In order to apply the previous proposition we will need the following result:

PROPOSITION 4. --- Let K be an ordered field and let \(f \in \mathbf{K}[\mathbf{X}]\) . There exists an interval in K, in the complement of which \(f\) takes the same sign as its highest degree term.

We can immediately reduce ourselves to the case of a monic polynomial; then one can write \(f(x) = x^{n}(1 + a_{1}x^{-1} + \cdots + a_{n}x^{-n})\) for \(x \neq 0\) . Let

\[
\mathbf {M} = \sup \left(1, | a _ {1} | + \dots + | a _ {n} |\right).
\]

For \(|x| > M\) we have \(1 + a_{1}x^{-1} + \cdots + a_{n}x^{-n} > 0\) , which completes the proof of the proposition.

COROLLARY 1. --- Every extension of an ordered field of odd finite degree admits the structure of an ordered extension.

Such an extension, being monogenous (V, p.~40, Th. 1), is isomorphic to K{[}X{]}/(f), where f is an irreducible polynomial of odd degree. Then it is enough to show that f changes sign in K (Prop. 3), which follows immediately from Prop. 4.

COROLLARY 2. --- If a is a positive element of an ordered field K, then every splitting field E of the polynomial \(X^{2}-a\) admits the structure of an ordered extension of K.

The result is trivial if a is a square in K. Otherwise the polynomial \(f(\mathbf{X}) = \mathbf{X}^{2} - a\) is irreducible and changes sign, since \(f(0) < 0\) and \(f(x)\) has the same sign as \(x^{2}\) , so positive, for x in the complement of some interval of K. We can now complete the proof by applying Prop. 3.

Remark. --- When the ordered field K contains the « square roots » of a positive element \(a\) of K (roots of the polynomial \(\mathbf{X}^2 - a\) ) then the notation \(\sqrt{a}\) is generally reserved for the positive square root. If K does not contain the square roots \(b\) and \(-b\) of \(a\) in the field E, then the latter can be made an ordered extension of K in two ways, each induced from the other via the K-automorphism which sends \(b\) to \(-b\) ; the choice of one of these orderings determines \(\sqrt{a}\) : it is whichever of the elements \(b\) and \(-b\) is positive.

If \(\pmb{a}\) and \(a'\) are two positive elements of \(\mathbf{K}\) , whose square roots are in \(\mathbf{K}\) , then \(\sqrt{aa'} = \sqrt{a}\sqrt{a'}\) , which follows from the definition of \(\sqrt{a}\) and the sign rule.

